\chapter{Math}
\section{Mathematical Preliminaries}
%====================
\subsection{Continuity of Derivatives}
Several results will make use
of the {\bf Schwarz theorem}, which 
states that if function $f : \real^2 \mapsto \real$ has continuous 
second partial derivatives at a given point $\vp \in \real^2 = 
\left\{p_1, p_2\right\}^T$,\footnote{This relationship can be 
extended to general dimensionality of 
$f : \real^n \mapsto \real$, but this higher dimensionality is not used
in the development that follows.}
then 
\be
 \dst\frac{\partial^2 f(\vp)}{\partial x_i \; \partial x_k} = 
 \dst\frac{\partial^2 f(\vp)}{\partial x_k \; \partial x_i}
 \hspace{1.0cm}\mbox{\bf (Commutative)}
 \label{eq:schwarz}
\ee
The Schwarz theorem establishes that the result of differentiation
of a function of two variables is independent of the order of 
that differentiation.


\subsection{Vector Spaces}
As defined by Gurtin,\footnote{Gurtin ME. An introduction to 
continuum mechanics. Academic press; 1982, Vol.~158, Jan 12.}
\begin{itemize}
\item \lin =  the set of all tensors
\item \linplus =  the set of all tensors $\tSbold$ with $\det \tSbold > 0$
\item \sym =  set of all symmetric tensors
\item \skw =  set of all skew tensors
\item \psym =  set of all positive definite, symmetric tensors
\item \orth =  set of all orthogonal tensors
\item \orthplus =  set of all rotations  
\end{itemize}



As defined by Curnier,\footnote{Curnier A. Computational methods 
in solid mechanics. Springer Science \& Business Media; 2012 Dec 6.}
for all second-order tensors $\tSbold \in$~\lin, 
the {\bf tracor} $\mathbb{T}$ and {\bf deviator} $\mathbb{D}$ 
are fourth-order tensors are defined as
\be
 \mathbb{T} \defe \frac{1}{3} \delta_{ij} \delta_{kl} \; 
 \veeihat \otimes
 \veejhat \otimes 
 \veekhat \otimes
 \veelhat 
\ee
such that 
\be
 \mathbb{T} \; \vddiv \; \tSbold = \tSbold_{\smVOL} 
 \defe \frac{1}{3} \trace(\tSbold)\vid, 
\ee
and
\be
 \mathbb{D} \defe 
 \left(
 \frac{1}{2} \delta_{ik} \delta_{jl} +
 \frac{1}{2} \delta_{il} \delta_{jk} - 
 \frac{1}{3} \delta_{ij} \delta_{kl} \; 
 \right)
 \veeihat \otimes
 \veejhat \otimes 
 \veekhat \otimes
 \veelhat 
\ee
such that
\be
 \mathbb{D} \; \vddiv \; \tSbold = \tSbold_{\smDEV}
 \defe \tSbold - \frac{1}{3} \trace(\tSbold)\vid.
\ee
Note, there are two other fourth-order tensors often used, 
the {\bf symmetrizer} $\mathbb{S}$ and the {\bf skewer} $\mathbb{W}$
defined, respectively, as 
\be
 \mathbb{S} \defe 
 \frac{1}{2}
 \left(
 \delta_{ik} \delta_{jl} +
 \delta_{il} \delta_{jk} 
 \right)
 \veeihat \otimes
 \veejhat \otimes 
 \veekhat \otimes
 \veelhat,
\ee
and
\be
 \mathbb{W} \defe 
 \frac{1}{2}
 \left(
 \delta_{ik} \delta_{jl} -
 \delta_{il} \delta_{jk} 
 \right)
 \veeihat \otimes
 \veejhat \otimes 
 \veekhat \otimes
 \veelhat.
\ee

\subsection{Matrix Operations}

\subsubsection{Symmetry Operators}
Let $\tAbold$ be a matrix of dimension $n \times n$, 
with components $a_{ij}$, 
$i=1\ldots n$, $j=1\ldots n$.  
The symmetry operator \sym($\cdot$), on any square square 
matrix $\tAbold$, is defined
\be
\mbox{\sym}(\tAbold) \defe \half (\tAbold + \tAbold\transpose).
\ee
The skew-symmetry operator \skw($\cdot$), on any square square 
matrix $\tAbold$, is defined
\be
\mbox{\skw}(\tAbold) \defe \half (\tAbold - \tAbold\transpose).
\ee
Then {\em any} matrix $\tAbold$ can be decomposed into symmetric and
skew-symmetric parts through an additive decomposition as
\be
\tAbold = \underset{\mbox{\sym}}{\underbrace{\dst\half\left(\tAbold 
+ \tAbold\transpose\right)}}
      \;\;  + \;\; \underset{\mbox{\skw}}{\underbrace{\dst\half\left(\tAbold 
- \tAbold\transpose\right)}} = 
\mbox{\sym}(\tAbold) \; + \mbox{\skw}(\tAbold).
\ee
Let the symmetric, square matrix $\tMbold = \mbox{\sym}(\tAbold)$ and
the skew-symmetric, square matrix $\tWbold = \mbox{\skw}(\tAbold)$.  Then
the matrix $\tAbold$ can be written as
\be
\tAbold = \tMbold + \tWbold.
\ee
The symmetric matrix $\tMbold$ has the property
\be
\tMbold = \tMbold\transpose 
\;\;\Longleftrightarrow\;\;
m_{ij} = m_{ji}.
\ee
Explicitly, for $\tMbold$ as a (3 $\times$ 3) matrix
\be
\left[
    \begin{tabular}{ccc}
    $m_{11}$ & $m_{12}$ & $m_{13}$ \\
    $m_{21}$ & $m_{22}$ & $m_{23}$ \\
    $m_{31}$ & $m_{32}$ & $m_{33}$
    \end{tabular}
\right]
= 
\left[
    \begin{tabular}{ccc}
    $m_{11}$ & $m_{21}$ & $m_{31}$ \\
    $m_{12}$ & $m_{22}$ & $m_{32}$ \\
    $m_{13}$ & $m_{23}$ & $m_{33}$
    \end{tabular}
\right],
\ee
and of the nine components of $\tMbold$, only six are independent, 
\be
\tMbold = \left[
    \begin{tabular}{ccc}
    $m_{11}$ & $m_{12}$ & $m_{13}$ \\
    $m_{12}$ & $m_{22}$ & $m_{23}$ \\
    $m_{13}$ & $m_{23}$ & $m_{33}$
    \end{tabular}
\right]
= 
\left[
    \begin{tabular}{ccc}
    $m_{11}$ & $m_{12}$ & $m_{13}$ \\
             & $m_{22}$ & $m_{23}$ \\
    $sym.$   &          & $m_{33}$
    \end{tabular}
\right].
\ee
The skew-symmetric matrix $\tWbold$ has the property
\be
\tWbold = -\tWbold\transpose 
\;\;\Longleftrightarrow\;\; w_{ij} = -w_{ji}.
\ee
Explicitly, for $\tWbold$ as a (3 $\times$ 3) matrix
\be
\tWbold
= 
\left[
    \begin{tabular}{ccc}
    $0$       & $-w_{12}$  & $w_{13}$ \\
    $w_{12}$ & $0$       & $-w_{23}$ \\
    $-w_{13}$ & $w_{23}$ & $0$
    \end{tabular}
\right],
\ee
and of the nine components of $\tWbold$, only three are independent.

\begin{Hremark}
In some references, skew-symmetric is called {\em antisymmetric}.  Here, 
    we will use {\em skew-symmetric}, or sometimes just {\em skew} as an
    abbreviated form.
\end{Hremark}


\begin{Hremark}
From spectral theory, the non-zero eigenvalues exist in pairs as purely 
    imaginary numbers $0 \pm \lambda_k i$, $\lambda \in \real$, with
    \begin{itemize}
        \item $k=1\ldots n$ for $\dim(\tWbold) = (n, n)$ and $n$ even and 
        \item $k=1\ldots(n - 1)/2$ for $\dim(\tWbold) = (n, n)$ and $n$ odd.
    \end{itemize}
    For even $n$ dimensions, there are no zero eigenvalues.  
    For odd $n$ dimensions, there is
    exactly one zero eigenvalue.
In the explicit example for $\dim(\tWbold) = (3,3)$, 
\be 
    \ve{\Lambda} = \mbox{\eig}(\tWbold) 
= 
\left[
    \begin{tabular}{ccc}
    $0$ & $\lambda$ & $0$ \\
    $-\lambda$ & $0$ & $0$ \\
    $0$ & $0$ & $0$
    \end{tabular}
\right].
\ee
\end{Hremark}

\begin{Hremark}
For the skew-symmetric matrix $\tWbold$, the tensor $\tZbold = \tWbold^2$ is 
negative semi-definite,
\be
    \vv\transpose \tZbold \vv \leq 0 \;\; \forall \mbox{ vectors } \vv \in \real^n.
\ee
\end{Hremark}

\begin{Hresult}
All $n \times n$ skew-symmetric matrices of odd dimension are singular.
\end{Hresult}

\begin{Hproof}
For skew-symmetric matrices, 
\be
\tWbold = -\tWbold\transpose.
\ee
From the properties of determinants
\be
\det(\tWbold) = \det(-\tWbold\transpose) = (-1)^n \det(\tWbold),
\ee
which, since $n$ is odd, becomes,
\be
\det(\tWbold) = - \det(\tWbold) = 0,
\ee
\end{Hproof}

\subsubsection{Cross-Product Operator}
One special form of a skew-symmetry operator is a 
{\bf cross-product operator}, defined as
\be
\label{eq:cross_operator}
[\; \va \;]_{\times} \defe
\left[
\begin{tabular}{ccc}
0 & $-a_3$ & $a_2$ \\
$a_3$ & 0 & $-a_1$ \\ 
$-a_2$ & $a_1$ & 0  
\end{tabular}
\right].
\ee
Thus $\va \times \vb = [ \; \va \;]_{\times} \vb$.



\subsection{Extrema}

\subsubsection{Arg Max}

The {\bf arguments of the maxima}, also known as 
{\bf arg max} or {\bf argmax}, at the 
values of the domain $x$ that maximize the range $f(x)$.  

\begin{Hexample}{
Defined over the domain $x \in [0, 2\pi]$, 
the function $f(x) = \cos(x) $ has  
maximum values of $1$ at $x=0$ and $x=2\pi$.  This notion is
written with the $\max$ as
\be
\underset{x\in[0,2\pi]}{\max} \; f(x) = 
\underset{x\in[0,2\pi]}{\max} \cos(x) = 1.
\ee
The {\bf arg max} is then 
\be
\underset{x\in[0,2\pi]}{\arg \max} \; f(x) = 
\underset{x\in[0,2\pi]}{\arg \max} \; \cos(x) = \{0, 2\pi\}.
\ee
More generally, for $x$ defined over the entire real number line,
\be
\underset{x\in\real}{\arg \max} \; f(x) = 
\underset{x\in\real}{\arg \max} \; \cos(x) = \{2\pi \; k\} \mbox{ for } k \in \integer.
\ee
}\end{Hexample}
  
